%% sections/04-safety.tex
%% Prerequisites & Safety Notices
%% IEC/IEEE 82079-1:2019 §7 — Safety Information
%% ANSI Z535.6 — Signal word conventions
%% ─────────────────────────────────────────────────

\section{Prerequisites and Safety Notices}
\label{sec:safety}

%%
%% WRITING GUIDE — SAFETY SECTION:
%%
%%   Per IEC/IEEE 82079-1, safety information MUST appear BEFORE the
%%   procedure, not embedded in steps.  Use DANGER for life-threatening
%%   hazards, WARNING for injury risk, CAUTION for equipment damage,
%%   NOTE for important operational information.
%%
%%   Subsections:
%%     4.1  Hardware prerequisites
%%     4.2  Software prerequisites
%%     4.3  Personnel prerequisites (training/competency)
%%     4.4  Safety notices
%%

\subsection{Hardware Prerequisites}
\label{sec:safety:hardware}

Confirm the following hardware is available and in working condition before
starting \cref{sec:setup}:

\begin{itemize}[noitemsep, topsep=4pt]
  \item[$\square$] PBA GS-Series gantry system, fully assembled and powered
  \item[$\square$] ACS SPiiPlus controller, firmware version $\geq$ 2.85
  \item[$\square$] Analogue load cell (Kistler 9317C or equivalent), factory
                   calibrated
  \item[$\square$] Signal conditioner / amplifier (±\SI{10}{\volt} output range)
  \item[$\square$] Shielded analogue cable, length $\leq$ \SI{3}{\metre}
  \item[$\square$] Personal computer with Ethernet connection to the ACS
                   controller
  \item[$\square$] Multimeter (for sensor output verification)
\end{itemize}


\subsection{Software Prerequisites}
\label{sec:safety:software}

\begin{itemize}[noitemsep, topsep=4pt]
  \item[$\square$] ACS SPiiPlus MMI application, version $\geq$ 2.85
                   (available from ACS Motion Control portal)
  \item[$\square$] Python 3.10 or newer
  \item[$\square$] Python libraries: \code{acspy}, \code{numpy},
                   \code{matplotlib} (install via \code{pip} — see
                   \cref{app:python-env})
  \item[$\square$] PBA configuration file
                   \filepath{GS500\_ForceCtrl\_baseline.cfg} (request from
                   PBA Systems support)
\end{itemize}


\subsection{Personnel Prerequisites}
\label{sec:safety:personnel}

This procedure should only be performed by personnel who have completed:

\begin{itemize}[noitemsep]
  \item PBA Systems GS-Series operator training (or equivalent)
  \item ACS SPiiPlus basic commissioning training
  \item Site-specific electrical safety induction
\end{itemize}

\begin{tipbox}
  University collaborators using this note as part of the PBA–University of
  Moratuwa research programme must also complete the PBA induction checklist
  (form PBA-HSE-001) before entering the lab.
\end{tipbox}


\subsection{Safety Notices}
\label{sec:safety:notices}

Read all notices below before proceeding. They remain applicable throughout
the entire procedure.

\bigskip

\begin{dangerbox}
  The Z-axis motor exerts forces up to \SI{500}{\newton}.  Never place
  fingers or hands under the Z-axis head while the system is powered and
  servos are enabled.  Always use the Emergency Stop (E-Stop) when
  adjusting hardware with the machine energised.
\end{dangerbox}

\bigskip

\begin{warningbox}
  Force control mode removes the position-based end-of-travel protection.
  An incorrectly tuned force loop or wrong force setpoint can drive the
  axis into hard mechanical contact at high current, damaging the machine,
  the workpiece, or the sensor.  Set a conservative force limit before
  enabling the force loop for the first time (see \cref{sec:procedure},
  Step~4).
\end{warningbox}

\bigskip

\begin{cautionbox}
  The load cell cable shield must be grounded at the signal conditioner end
  only (single-point grounding).  Grounding at both ends introduces a
  ground loop that causes \SI{50}{\hertz} / \SI{60}{\hertz} noise on the
  force signal, leading to loop instability.
\end{cautionbox}

\bigskip

\begin{notebox}
  All parameter changes made during this procedure are non-persistent by
  default on the ACS SPiiPlus controller.  Use the \menu{File > Save
  Configuration} command in the MMI after confirming stable operation to
  write parameters to non-volatile memory.
\end{notebox}
